Gracias a este análisis empírico se demuestra que el rendimiento práctico de un algoritmo depende tanto de su complejidad teórica como de las constantes ocultas de su implementación. En el problema de ordenamiento, \textit{Quick Sort} y \textit{Standard Sort} demostraron ser las opciones más eficientes en tiempo y memoria gracias a sus particiones \textit{in-place}, mientras que \textit{Merge Sort} validó su robustez asintótica sacrificando un espacio de memoria adicional $\Theta(n)$. Por otro lado, en la multiplicación de matrices, los datos probaron de manera contundente el impacto negativo de las constantes ocultas en métodos recursivos, a pesar de tener una mejor cota teórica, \textit{Strassen} fue ampliamente superado por el enfoque \textit{Naive} para dimensiones de hasta $N=2^{10}$. Esto se debió a su excesiva sobrecarga de memoria dinámica y a la excelente optimización de localidad de referencia (caché) implementada en los bucles de \textit{Naive}. Finalmente, las pruebas confirmaron que ambos algoritmos de matrices son insensibles al grado de dispersión de los datos, ejecutando su respectiva complejidad asintótica independientemente del contenido de la estructura.